\chapter{时间旅行的概率动力学研究}

本章内容是模板示例，用于展示理论基础、概率模型、数学分析和中文定义/定理环境的组织方式。本文中的时间分支模型是形式化算例，不构成关于现实时间旅行可行性的证据。

\section{理论基础}

\subsection{相对论框架}

示例内容从参考系、时间参数和状态转移三个方面说明模型的基本假设。相对论中的时空结构决定事件之间的因果关系；若形式模型允许沿闭合类时曲线返回较早的时刻，还需要增加一致性条件，防止状态沿一周演化后与初始条件矛盾。这里仅把这些条件作为离散状态模型的约束来源，不尝试推导新的时空解。

\subsection{量子效应}

若研究涉及离散状态或不确定性，可在本节说明状态空间、转移概率以及观测量的定义。为展示中文学位论文中的定理环境，本章采用无边框的定义、定理与证明格式。

\begin{defn}[可逆状态转移]
若状态转移算子 $\bm{A}$ 可逆，则称该离散模型在给定状态空间上可逆。
\end{defn}

\begin{thm}[状态转移的可逆性]
若 $\bm{A}$ 可逆且 $\bm{x}_{k+1}=\bm{A}\bm{x}_k$，则由终态 $\bm{x}_{k+1}$ 可以唯一恢复 $\bm{x}_k$。
\end{thm}

\begin{proof}
左乘 $\bm{A}^{-1}$ 得 $\bm{x}_k=\bm{A}^{-1}\bm{x}_{k+1}$。由于 $\bm{A}^{-1}$ 存在，给定终态后恢复结果唯一。
\end{proof}

上述定义和定理的编号会随章节自动重置；真实论文可继续增加引理、推论和定理证明，而不需要手工维护编号。

为了展示状态关系的几种表达方式，图~\ref{fig:ch3-relations}并列给出交换关系、范畴关系、离散状态机和层级状态树。交换图强调不同路径之间的一致性，范畴图表示对象与态射之间的映射，状态机描述状态转移，层级树则呈现对象的分类结构。正式研究应在图前定义节点集合和映射，在图后说明图中的关系如何对应式~\eqref{eq:state-update}。
\begin{figure}[htbp]
  \centering
  % The four panels share the content and two-by-two layout of the source figure.
  \begin{subfigure}[t]{0.49\textwidth}
    \centering
    \resizebox{\linewidth}{!}{%
      \begin{tikzpicture}[>=Stealth, every node/.style={font=\scriptsize}]
        \useasboundingbox (0,-0.8) rectangle (6,2.4);
        \node (top-left) at (2.05,1.911)
          {$\tau^{*}\!\left(AB^{n}\dfrac{R}{R^{\times}}\right)$};
        \node (top-right) at (4.75,1.911)
          {$B^{n}\dfrac{R}{R^{\times}}$};
        \node (left) at (0.55,0.367) {$X$};
        \node (middle) at (2.8,0.367) {$X$};
        \node (right) at (5.0,0.367) {$\boldsymbol{B}R^{\times}$};
        \draw[->] (top-left) -- (top-right);
        \draw[->] (top-left.south) -- (middle.north);
        \draw[->] (top-right.south) -- (right.north);
        \draw[dashed,->] (left.north east) -- (top-left.south west);
        \draw[double,double distance=1.2pt] (left.east) -- (middle.west);
        \draw[->] (middle) -- node[above] {$\tau$} (right);
        \draw[->,blue,thick] (right.south) -- (5.0,-0.578) --
          (0.55,-0.578) -- (left.south);
        \node[blue,fill=white,inner sep=1pt] at (2.78,-0.578) {$\Gamma_R^{*}$};
      \end{tikzpicture}%
    }
    \caption{交换关系}
  \end{subfigure}\hfill
  \begin{subfigure}[t]{0.49\textwidth}
    \centering
    \resizebox{\linewidth}{!}{%
      \begin{tikzpicture}[>=Stealth, every node/.style={font=\scriptsize}]
        \useasboundingbox (0,-0.8) rectangle (6,2.4);
        \node (id) at (0.45,1.806) {$\mathrm{id}_S\in$};
        \node (hom-top) at (1.95,1.806) {$\operatorname{Hom}_{\mathcal C}(S,S)$};
        \node (a-top) at (4.00,1.806) {$A(S)$};
        \node (u) at (5.55,1.806) {$\ni u_S$};
        \node (f) at (0.45,0.263) {$f\in$};
        \node (hom-bottom) at (1.95,0.263) {$\operatorname{Hom}_{\mathcal C}(T,S)$};
        \node (a-bottom) at (4.00,0.263) {$A(T)$};
        \node (phi) at (5.55,0.263) {$\ni\phi_T(f)$};
        \draw[->] (hom-top) -- node[above] {$\phi_S$} (a-top);
        \draw[->] (hom-bottom) -- node[above] {$\phi_T$} (a-bottom);
        \draw[->] (id) -- (f);
        \draw[->] (hom-top) -- node[right] {$f^{*}$} (hom-bottom);
        \draw[->] (a-top) -- node[right] {$A(f)$} (a-bottom);
        \draw[->] (u) -- (phi);
      \end{tikzpicture}%
    }
    \caption{范畴关系}
  \end{subfigure}

  \par\smallskip

  \begin{subfigure}[t]{0.49\textwidth}
    \centering
    \resizebox{\linewidth}{!}{%
      \begin{tikzpicture}[>=Stealth, every node/.style={font=\scriptsize},
        state/.style={draw,circle,minimum size=12mm,inner sep=1pt,fill=blue!12}]
        \useasboundingbox (0,-0.8) rectangle (6,2.4);
        \node[state] (reading) at (1.45,0.473) {reading};
        \node[state] (eof) at (3.35,0.473) {eof};
        \node[state,double,double distance=2pt] (closed) at (5.2,0.473) {closed};
        \draw[->] (0.05,0.473) -- node[above] {open(path)} (reading.west);
        \draw[->] (reading.east) -- node[above] {read()} (eof.west);
        \draw[->] (eof.east) -- node[above] {close()} (closed.west);
        \draw[->] (reading) to[loop above] node[above=2pt] {read()} (reading);
        \draw[->] (reading.south east) to[bend right=32]
          node[below] {close()} (closed.south west);
      \end{tikzpicture}%
    }
    \caption{离散状态机}
  \end{subfigure}\hfill
  \begin{subfigure}[t]{0.49\textwidth}
    \centering
    \resizebox{\linewidth}{!}{%
      \begin{tikzpicture}[>=Stealth, every node/.style={font=\scriptsize},
        tree-node/.style={draw=gray!55,fill=white,align=center,font=\tiny,
          text width=13mm,minimum height=6mm,inner sep=1.5pt},
        tree-node-wide/.style={tree-node,text width=18mm}]
        \useasboundingbox (0,-0.8) rectangle (6,2.4);
        \node[tree-node] (root) at (3.0,2.047) {Entrate\\pubblicitarie};
        \node[tree-node] (impr) at (1.55,0.892) {Numero di\\impression};
        \node[tree-node-wide] (gain) at (4.45,0.892) {Guadagno medio\\per impression};
        \node[tree-node] (traffic) at (0.55,-0.294) {Traffico\\sul sito};
        \node[tree-node] (ratio) at (2.25,-0.294) {Impression/\\visitatori};
        \node[tree-node] (sales) at (3.75,-0.294) {Modalità\\di vendita};
        \node[tree-node] (placement) at (5.45,-0.294) {Tipo di\\posi\-zionamento};
        \draw[->] (root) |- ($(root)!0.5!(impr)$) -| (impr);
        \draw[->] (root) |- ($(root)!0.5!(gain)$) -| (gain);
        \draw[->] (impr) |- ($(impr)!0.5!(traffic)$) -| (traffic);
        \draw[->] (impr) |- ($(impr)!0.5!(ratio)$) -| (ratio);
        \draw[->] (gain) |- ($(gain)!0.5!(sales)$) -| (sales);
        \draw[->] (gain) |- ($(gain)!0.5!(placement)$) -| (placement);
      \end{tikzpicture}%
    }
    \caption{层级状态树}
  \end{subfigure}
  \caption{状态关系图组：交换图、范畴图、状态机与层级树}
  \label{fig:ch3-relations}
\end{figure}

\section{概率模型}

\subsection{模型参数与边界条件}

离散时间模型由状态向量、控制量、过程噪声和观测噪声构成。初始分布、转移参数和边界条件应由研究对象及其数据来源确定；在本模板中，这些量均作为可替换的示例记号。

\subsection{行内、行间与分段公式}

短表达式适合放在正文中，例如状态转移的离散形式可写为 $\bm{x}_{k+1}=\bm{A}\bm{x}_k+\bm{w}_k$。当公式需要独立展示时，使用带编号的行间公式：
\begin{equation}
  \bm{x}_{k+1}=\bm{A}\bm{x}_k+\bm{B}\bm{u}_k+\bm{w}_k.
  \label{eq:state-update}
\end{equation}

分段定义可用 \texttt{cases} 环境表示。与普通文本一样，分段公式会自动继承模板的数学字号和章节编号规则。
\begin{equation}
  P(X=x_i)=
  \begin{cases}
    p_i, & x_i\in S,\\
    0, & x_i\notin S.
  \end{cases}
  \label{eq:piecewise-probability}
\end{equation}

对于同时给出观测方程和状态方程的场景，可用 \texttt{aligned} 环境将多行内容排成一个公式块：
\begin{equation}
  \begin{aligned}
    \bm{x}_{k+1}&=\bm{A}\bm{x}_k+\bm{B}\bm{u}_k+\bm{w}_k,\\
    \bm{y}_k&=\bm{C}\bm{x}_k+\bm{v}_k.
  \end{aligned}
  \label{eq:state-observation}
\end{equation}

公式标签也可以跨章节引用，例如式~\eqref{eq:state-update}。真实论文中建议为关键定义、定理和推导保留稳定标签，避免手工维护编号。

\subsection{模型验证与讨论}

为了验证状态转移模型，需要将预测状态与观测状态放在同一时间轴上比较。若观测误差协方差记为 $R$，则一个简单的标准化残差可以写成
\begin{equation}
  z_k=\frac{\bm{y}_k-\bm{C}\bm{x}_k}{\sqrt{R}}.
  \label{eq:standardized-residual}
\end{equation}

当 $z_k$ 长时间偏离零时，应优先检查观测矩阵、边界条件和时间同步，而不是直接增加模型参数。这个例子也说明，行间公式、解释文字和图表之间应形成连续的推理链；公式不能脱离上下文单独占据页面。

对于多个状态变量，可以进一步报告状态转移矩阵的谱半径、预测误差和收敛步数。正式论文中应在本节给出实验设置、重复次数和统计检验方法，并说明哪些结果来自真实数据，哪些结果只是数值模拟。

\section{数学分析基础}

\subsection{极限、连续与导数}

时间演化的局部性质由导数描述。把位移、速度和加速度写成对时间的各阶导数后，模型的状态量之间就建立了可微的依赖关系，后续的线性化和灵敏度分析都建立在这一步之上。
\begin{equation}
  \lim_{h\to0}\frac{f(x_0+h)-f(x_0)}{h}=f'(x_0),\qquad
  f'(x)=\frac{\mathrm{d}y}{\mathrm{d}x},\qquad
  f''(x)=\frac{\mathrm{d}^2y}{\mathrm{d}x^2},\qquad
  \dot{x}(t)=\frac{\mathrm{d}x}{\mathrm{d}t}.
  \label{eq:derivative}
\end{equation}

中值定理和泰勒展开把局部导数与整体变化联系起来。工程上常取二阶展开近似代价函数在最优解附近的形状，据此判断迭代收敛是线性还是二次的。
\begin{equation}
  f(x_0+h)=f(x_0)+f'(x_0)h+\frac{1}{2}f''(x_0)h^2+O(h^3).
  \label{eq:taylor}
\end{equation}
\begin{equation}
  f(b)-f(a)=f'(\xi)(b-a),\qquad \xi\in(a,b).
  \label{eq:mean-value}
\end{equation}

多尺度模型的渐近分析还要求处理不定式极限。下面两式分别给出小量展开和远场衰减的常用结果，正文中应说明展开成立的量级范围。
\begin{equation}
  \lim_{x\to0}\frac{\sin x}{x}=1,\qquad
  \lim_{x\to\infty}\frac{\ln x}{x^p}=0\quad(p>0),\qquad
  e^x=\sum_{k=0}^{\infty}\frac{x^k}{k!}.
  \label{eq:asymptotic-limits}
\end{equation}

\subsection{最优性条件}

无约束问题的一阶必要条件给出驻点，二阶条件区分极小点与鞍点。把梯度写成向量、把二阶导数写成 Hessian 矩阵后，两个条件可以用同一套记号表达。
\begin{equation}
  \min_{\bm{x}\in\mathbb{R}^n}f(\bm{x}),\qquad
  \nabla f(\bm{x}^*)=\bm{0},\qquad
  \bm{H}(\bm{x}^*)=
  \left(\frac{\partial^2f}{\partial x_i\partial x_j}\right)_{ij}\succ0.
  \label{eq:optimality}
\end{equation}

带约束的问题用拉格朗日函数处理：等式约束引入乘子，不等式约束引入非负乘子。KKT 条件把原始可行性与对偶可行性、互补松弛写在同一组方程里。
\begin{equation}
  \mathcal{L}(\bm{x},\bm{\lambda},\bm{\mu})
  =f(\bm{x})+\sum_{i=1}^{m}\lambda_i g_i(\bm{x})
  +\sum_{j=1}^{p}\mu_j h_j(\bm{x}).
  \label{eq:lagrangian}
\end{equation}
\begin{equation}
  \nabla f(\bm{x}^*)+\sum_{i=1}^{m}\lambda_i\nabla g_i(\bm{x}^*)
  +\sum_{j=1}^{p}\mu_j\nabla h_j(\bm{x}^*)=\bm{0},\qquad
  \mu_j\geq0,\qquad \mu_jh_j(\bm{x}^*)=0.
  \label{eq:kkt}
\end{equation}

当目标函数是凸函数、可行域是凸集时，KKT 条件不仅是必要条件，也是充分条件。这一性质决定了内点法和序列二次规划能否给出全局最优解，正文应明确所求解问题的凸性。
\begin{equation}
  f(\theta\bm{x}+(1-\theta)\bm{y})
  \leq\theta f(\bm{x})+(1-\theta)f(\bm{y}),\qquad \theta\in[0,1].
  \label{eq:convexity}
\end{equation}

\subsection{矩阵分解与谱}

对称正定矩阵的特征分解把二次型化为平方和，条件数衡量问题对扰动的敏感程度。谱半径同时决定迭代格式是否收敛，因此它既出现在数值分析里，也出现在稳定性讨论中。
\begin{equation}
  \bm{A}=\bm{Q}\bm{\Lambda}\bm{Q}^{\mathsf T},\qquad
  \bm{A}\bm{q}_i=\lambda_i\bm{q}_i,\qquad
  \bm{x}^{\mathsf T}\bm{A}\bm{x}
  =\sum_{i=1}^{n}\lambda_i(\bm{q}_i^{\mathsf T}\bm{x})^2.
  \label{eq:eigen}
\end{equation}

一般矩阵用奇异值分解。截断较小的奇异值就是在最小二乘意义下取得最优的低秩近似，这一点在模型降阶和参数矩阵压缩中直接使用。
\begin{equation}
  \bm{A}=\bm{U}\bm{\Sigma}\bm{V}^{\mathsf T}
  =\sum_{i=1}^{r}\sigma_i\bm{u}_i\bm{v}_i^{\mathsf T},\qquad
  \kappa(\bm{A})=\frac{\sigma_{\max}}{\sigma_{\min}}.
  \label{eq:svd}
\end{equation}

\subsection{不等式与误差估计}

推导误差上界时常用的几个不等式如下。它们给出的是最坏情况估计，实际误差通常更小，因此正文应说明估计的保守程度，而不应把上界当作实际误差报告。
\begin{equation}
  |\bm{x}^{\mathsf T}\bm{y}|
  \leq\|\bm{x}\|_2\|\bm{y}\|_2,\qquad
  \left(\sum_{i=1}^{n}a_ib_i\right)^2
  \leq\left(\sum_{i=1}^{n}a_i^2\right)
       \left(\sum_{i=1}^{n}b_i^2\right).
  \label{eq:cauchy-schwarz}
\end{equation}
\begin{equation}
  \varphi(\mathbb{E}[X])\leq\mathbb{E}[\varphi(X)]
  \quad(\varphi\text{ 为凸函数}),\qquad
  \|\bm{x}+\bm{y}\|\leq\|\bm{x}\|+\|\bm{y}\|.
  \label{eq:jensen}
\end{equation}

把上述工具组合起来，可以给出数值解与真解之差的先验估计。下式中右端的第一项来自截断误差，第二项来自数据误差，二者随网格加密的下降速度不同，因此需要分别报告。
\begin{equation}
  \|\bm{u}_h-\bm{u}\|\leq C_1h^p+C_2\delta,\qquad
  \delta=\frac{\|\bm{y}-\bm{y}_{\mathrm{true}}\|}
  {\|\bm{y}_{\mathrm{true}}\|}.
  \label{eq:error-bound}
\end{equation}

\section{本章小结}

本章示例给出了理论基础、概率模型、参数设置以及中文定理环境的组织方式。定理环境的视觉样式属于内容层扩展，不会修改模板的页眉、正文间距或章节编号规则。
